% from ega1-9-fr.tex
\oldpage[I]{177}
est un second sous-préschéma fermé de $Y$, défini par un faisceau
quasi-cohérent d'idéaux $\mathscr{J}'$ de $\mathscr{O}_Y$, et tel que
l'injection $j':X''\to Y$ majore $f$, on doit tout d'abord avoir
$X''\supset\psi(X)$, donc $X'\subset X''$ puisque $X''$ est fermé. En outre,
pour tout $y\in X''$, $\theta$ doit se factoriser en
$\mathscr{O}_y\to\mathscr{O}_y/\mathscr{J}'_y
\to(\psi_*(\mathscr{O}_X))_y$, ce qui par définition entraîne
$\mathscr{J}'_y\subset\mathscr{J}_y$, et par suite $X'$ est un
sous-préschéma fermé de $X''$ (\hyperref[I.4.1.10-fr]{4.1.10}).
\end{proof}

\begin{definition}[9.5.3]
\label{I.9.5.3-fr}
Lorsqu'il existe un plus petit sous-préschéma fermé $Y'$ de $Y$ tel que
l'injection canonique $j:Y'\to Y$ majore $f$, on dit que $Y'$ est le
préschéma image fermée de $X$ par le morphisme $f$.
\end{definition}

\begin{proposition}[9.5.4]
\label{I.9.5.4-fr}
Si $f_*(\mathscr{O}_X)$ est un $\mathscr{O}_Y$-Module quasi-cohérent,
l'espace sous-jacent de l'image fermée de $X$ par $f$ est l'adhérence
$\overline{f(X)}$ dans $Y$.
\end{proposition}

\begin{proof}
Comme le support de $f_*(\mathscr{O}_X)$ est contenu dans
$\overline{f(X)}$, on a (avec les notations de
(\hyperref[I.9.5.2-fr]{9.5.2})) $\mathscr{J}_y=\mathscr{O}_y$ pour
$y\notin\overline{f(X)}$, donc le support de
$\mathscr{O}_Y/\mathscr{J}$ est contenu dans $\overline{f(X)}$. En outre, ce
support est fermé et contient $f(X)$~: en effet, si $y\in f(X)$, l'élément
unité de l'anneau $(\psi_*(\mathscr{O}_X))_y$ n'est pas nul, étant le germe
en $y$ de la section
\[
  1\in\Gamma(X,\mathscr{O}_X)=\Gamma(Y,\psi_*(\mathscr{O}_X))~;
\]
comme il est l'image par $\theta$ de l'élément unité de $\mathscr{O}_y$, ce
dernier n'appartient pas à $\mathscr{J}_y$, donc
$\mathscr{O}_y/\mathscr{J}_y\neq0$~; ceci achève la démonstration.
\end{proof}

\begin{proposition}[9.5.5]
\label{I.9.5.5-fr}
\emph{(transitivité des images fermées).} Soient $f:X\to Y$ et $g:Y\to Z$
deux morphismes de préschémas~; on suppose que l'image fermée $Y'$ de $X$
par $f$ existe, et que, si $g'$ est la restriction de $g$ à $Y'$, l'image
fermée $Z'$ de $Y'$ par $g'$ existe. Alors l'image fermée de $X$ par
$g\circ f$ existe et est égale à $Z'$.
\end{proposition}

\begin{proof}
Il suffit (\hyperref[I.9.5.1-fr]{9.5.1}) de montrer que $Z'$ est le plus
petit sous-préschéma fermé $Z_1$ de $Z$ tel que le sous-préschéma fermé
$(g\circ f)^{-1}(Z_1)$ de $X$ (égal à $f^{-1}(g^{-1}(Z_1))$ par
(\hyperref[I.4.4.2-fr]{4.4.2})) soit égal à $X$~; il revient au même de dire
que $Z'$ est le plus petit sous-préschéma fermé de $Z$ tel que $f$ soit
majoré par l'injection $g^{-1}(Z_1)\to Y$
(\hyperref[I.4.4.1-fr]{4.4.1}). Or, en vertu de l'existence de l'image
fermée $Y'$, tout $Z_1$ ayant cette propriété est tel que $g^{-1}(Z_1)$
majore $Y'$, ce qui équivaut à dire que
$j^{-1}(g^{-1}(Z_1))=g'^{-1}(Z_1)=Y'$, en désignant par $j$ l'injection
$Y'\to Y$. Par définition de $Z'$, on en conclut bien que $Z'$ est le plus
petit sous-préschéma fermé de $Z$ vérifiant la condition précédente.
\end{proof}

\begin{corollary}[9.5.6]
\label{I.9.5.6-fr}
Soit $f:X\to Y$ un $S$-morphisme tel que $Y$ soit l'image fermée de $X$ par
$f$. Soit $Z$ un $S$-schéma~; si deux $S$-morphismes $g_1$, $g_2$ de $Y$
dans $Z$ sont tels que $g_1\circ f=g_2\circ f$, alors $g_1=g_2$.
\end{corollary}

\begin{proof}
Soit $h=(g_1,g_2)_S:Y\to Z\times_S Z$~; comme la diagonale
$T=\Delta_Z(Z)$ est un sous-préschéma fermé de $Z\times_S Z$,
$Y'=h^{-1}(T)$ est un sous-préschéma fermé de $Y$
(\hyperref[I.4.4.1-fr]{4.4.1}). Posons $u=g_1\circ f=g_2\circ f$~; on a
alors par définition du produit $h'=h\circ f=(u,u)_S$, donc
$h\circ f=\Delta_Z\circ u$~; comme $\Delta_Z^{-1}(T)=Z$, on a
$h'^{-1}(T)=u^{-1}(Z)=X$, donc $f^{-1}(Y')=X$. On en conclut
(\hyperref[I.4.4.1-fr]{4.4.1}) que l'injection canonique $Y'\to Y$ majore
$f$, donc $Y'=Y$ par hypothèse~; par suite
(\hyperref[I.4.4.1-fr]{4.4.1}), $h$ se factorise en $\Delta_Z\circ v$, où
$v$ est un morphisme $Y\to Z$, ce qui entraîne $g_1=g_2=v$.
\end{proof}

\begin{remark}[9.5.7]
\label{I.9.5.7-fr}
Si $X$ et $Y$ sont des $S$-schémas, la prop.
(\hyperref[I.9.5.6-fr]{9.5.6}) signifie que
